\scanpage{167}
% archive foliation: 166
\inserted{Supposons \deleted{$\unsure{S}$} $S$ loc. noeth \uline{ou} $\pi$ artinien [condition des chaînes descendantes].}

(ii) Tout élément $c\in\Phi(\pi)$ est induit par un élément \unsure{minimal}, $c'\in\Phi(\pi')$.

\margin{\unsure{pour que} $\unsure{\Phi(\pi)}$ \unsure{ait un} \uline{\unsure{élémt}}}

En effet, si $c$ correspond à $P,\xi$, supposons \inserted{\unsure{corresp. d'éléments}} qu'il puisse \deleted{\unsure{être}} \unsure{réduire} à $P',\xi'$ \inserted{[groupe $\pi'\subset\pi$]} et $(P'',\xi'')$ [groupe $\pi''\subset\pi$].
Considérons
\[
P/_{\pi'\cap\pi''}\longrightarrow P/_{\pi'}\times P/_{\pi''}
\]
\margin{$\pi/_{\pi'\cap\pi''}\to\pi/_{\pi'}\times\pi/_{\pi''}$}
c'est une imm. ouverte et fermée, et d'autre part on a une section du $2^\circ$ \unsure{mbr} qui \deleted{est} au pt $(\xi',\xi'')$ est dans l'image du premier, elle y est donc partout.

Donc on peut restreindre le groupe structural à $\pi'\cap\pi''$. Donc les groupes auxquels on peut restreindre forment un ens filtrant décroissant, $\pi_i,\ldots$, donc $P_i\subset P$ forment un syst. projectif.

Si \deleted{\unsure{les groupes}} $\pi$ est artinien, \unsure{c'est fini}. Le système des $P_i$ est \deleted{\unsure{les uns inclus dans les autres} $P_i$ \unsure{est}} stationnaire. Si $S$ est loc. noeth., \deleted{\unsure{on}} \deleted{fini. Donc on peut former le préschéma $\varprojlim_i P_i$ qui est \unsure{encore fibré sur} $P$, \deleted{\unsure{le groupe est}} \inserted{et qui comme sous-pré\-schéma est} ouvert et fermé, et est un fibré}
